\paragrafos{Πρόσθεση και αφαίρεση ακεραίων: κέρδη, ζημιές και διαφορές}{Α.2.3, «Πρόσθεση και αφαίρεση ακεραίων»}

\begin{anagke}
Στο πούλμαν η παρέα παίζει ένα παιχνίδι με κάρτες: κάθε κάρτα γράφει έναν πόντο, θετικό για κέρδος και αρνητικό για ζημιά. Ο Βασίλης τραβάει $+5$ και μετά $-3$. «Πόσους πόντους έχω;» ρωτάει. Ο Νίκος υπολογίζει: «Πέντε μείον τρία κάνει δύο». Η Πολυξένη τραβάει $-4$ και μετά $-2$ και ο Στέλιος λέει «έξι», μα η Ρήνη τον ρωτάει: «Έξι κέρδος ή έξι ζημιά;». Ο Ηλίας κάνει μια πιο δύσκολη ερώτηση: «Αν τραβήξω $-7$ και μετά $+3$, είμαι ακόμα πίσω;».

Η Αγγελική το βλέπει στον άξονα: «Κάθε κάρτα είναι μια μετακίνηση. Το $+5$ σημαίνει «πέντε βήματα δεξιά», το $-3$ «τρία βήματα αριστερά». Ξεκινάω από το μηδέν, πηγαίνω πέντε δεξιά και μετά τρία αριστερά και καταλήγω στο $+2$». Για το $-4$ και το $-2$ είναι έξι βήματα αριστερά: $-6$. Και για τα $-7$ και $+3$ καταλήγουμε στο $-4$: ακόμα πίσω, αλλά λιγότερο. Ο Νίκος παρατηρεί και κάτι ωραίο: αν μια κάρτα $+3$ και μια $-3$ βρεθούν μαζί, «ακυρώνονται».

Το βράδυ ο μετεωρολόγος στην τηλεόραση λέει: «Η μέγιστη θερμοκρασία σήμερα είναι $+5$ βαθμοί και η ελάχιστη $-3$. Η διαφορά είναι $8$ βαθμοί». Ο Στέλιος κάνει τον λογαριασμό: «Πέντε μείον τρία κάνει δύο. Ο μετεωρολόγος κάνει λάθος». Η Πολυξένη τον διορθώνει: «Δεν αφαιρούμε το τρία από το πέντε. Αφαιρούμε το μείον τρία!». Η Αγγελική μετράει στον άξονα: από το $-3$ ως το $+5$ είναι τρία βήματα ως το μηδέν και πέντε ακόμη, οκτώ. Ο Ηλίας το γενικεύει: «Μήπως η αφαίρεση του $-3$ είναι το ίδιο με την πρόσθεση του $+3$;».

Σε αυτή την παράγραφο η παρέα ανακαλύπτει τους κανόνες της πρόσθεσης και της αφαίρεσης ακεραίων, μαθαίνει γιατί ισχύουν και πώς γίνονται σύντομοι οι υπολογισμοί όταν έχουμε πολλούς όρους και παρενθέσεις.
\end{anagke}

\begin{lexi}{πρόσθεση, άθροισμα, αφαίρεση, διαφορά}{πρόσθεση, άθροισμα, αφαίρεση, διαφορά}
Η \lexiword{πρόσθεση} είναι το «προσ-τίθημι», δηλαδή «βάζω επιπλέον». Οι \lexiword{προσθετέοι} είναι αυτά που προσθέτουμε και το \lexiword{άθροισμα} είναι αυτό που βγαίνει όταν τα «αθροίζουμε», τα βάζουμε σε ένα σωρό. Η \lexiword{αφαίρεση} είναι το «αφ-αιρώ», «βγάζω από κάτι». Η \lexiword{διαφορά} είναι αυτό που μένει και δείχνει «πόσο διαφέρουν» δύο αριθμοί. Στους ακέραιους, το «προσθέτω» μπορεί να σημαίνει και «ζημιώνω».
\end{lexi}

\begin{metafrasi}
$5$ βήματα δεξιά και $3$ αριστερά. & $(+5)+(-3)=+2$\\
$4$ βήματα αριστερά και $2$ αριστερά ακόμα. & $(-4)+(-2)=-6$\\
Κέρδος και ζημιά που ακυρώνονται. & $(+3)+(-3)=0$\\
Δεν αλλάζει η σειρά των προσθετέων. & $\alpha+\beta=\beta+\alpha$\\
Αφαίρεση σημαίνει πρόσθεση του αντίθετου. & $\alpha-\beta=\alpha+(-\beta)$\\
Η διαφορά θερμοκρασιών $+5$ και $-3$. & $(+5)-(-3)=8$\\
\end{metafrasi}

\tmima{Η θεωρία}

\tmimasmall{Πρόσθεση ακεραίων}

Θεωρούμε ένα ρομπότ που ξεκινά από την αρχή $\mathrm{Ο}$ του άξονα και κινείται μόνο πάνω στην ευθεία. Μια εντολή $+4$ το στέλνει $4$ μονάδες δεξιά και μια εντολή $-3$ το στέλνει $3$ μονάδες αριστερά. Αν δώσουμε δύο εντολές, η θέση που καταλήγει είναι το άθροισμα των δύο αριθμών.

\begin{center}
\begin{tikzpicture}[x=0.65cm]
  \draw[-{Stealth},thick] (-7.5,0)--(8.5,0);
  \foreach \j in {-7,...,8}{\draw (\j,0.1)--(\j,-0.1) node[below,font=\scriptsize]{$\j$};}
  \draw[thick] (0,0.17)--(0,-0.17);
  \draw[-{Stealth},gblue,thick] (0,0.45) to[bend left=25] node[above,font=\small]{$+4$} (4,0.45);
  \draw[-{Stealth},gwarmD,thick] (4,0.45) to[bend right=25] node[above,font=\small]{$-3$} (1,0.45);
  \fill[gblue] (1,0) circle (2.3pt);
\end{tikzpicture}
\end{center}

Το ρομπότ πήγε $4$ δεξιά και μετά $3$ αριστερά και έφτασε στο $+1$: $(+4)+(-3)=+1$. Αν κινηθεί $3$ αριστερά και μετά $1$ αριστερά, φτάνει στο $-4$: $(-3)+(-1)=-4$. Από αυτά βγαίνουν οι κανόνες:

\begin{kanonas}{πρόσθεση ακεραίων}
\begin{itemize}[leftmargin=1.2em,itemsep=0.3ex,topsep=2pt]
\item \textbf{Ομόσημοι:} προσθέτω τις απόλυτες τιμές και κρατάω το κοινό πρόσημο. $(-8)+(-5)=-(8+5)=-13$ και $(+9)+(+4)=+13$.
\item \textbf{Ετερόσημοι:} αφαιρώ τη μικρότερη απόλυτη τιμή από τη μεγαλύτερη και βάζω το πρόσημο του προσθετέου με τη μεγαλύτερη απόλυτη τιμή. $(+9)+(-4)=+5$ και $(-9)+(+4)=-5$.
\item \textbf{Αντίθετοι:} το άθροισμα είναι $0$. $(-6)+(+6)=0$.
\item Με το $0$: $\alpha+0=\alpha$.
\end{itemize}
\end{kanonas}

\begin{idiotita}{ιδιότητες της πρόσθεσης}
Για όλους τους ακέραιους $\alpha,\,\beta,\,\gamma$ ισχύουν: \textbf{αντιμεταθετική} $\alpha+\beta=\beta+\alpha$, \textbf{προσεταιριστική} $(\alpha+\beta)+\gamma=\alpha+(\beta+\gamma)$, $\alpha+0=\alpha$ και $\alpha+(-\alpha)=0$.
\end{idiotita}

Χάρη στις ιδιότητες, όταν έχω πολλούς προσθετέους μπορώ να τους ομαδοποιήσω όπως θέλω: προσθέτω χωριστά όλους τους θετικούς και χωριστά όλους τους αρνητικούς και μετά τα δύο αποτελέσματα. Έτσι $(-7)+12+(-5)+9+(-3)=(12+9)+\big((-7)+(-5)+(-3)\big)=21+(-15)=6$.

\tmimasmall{Αφαίρεση ακεραίων}

\begin{kanonas}{αφαίρεση ακεραίων}
Για να αφαιρέσω από τον $\alpha$ τον $\beta$, προσθέτω στον $\alpha$ τον \textbf{αντίθετο} του $\beta$: $$\alpha-\beta=\alpha+(-\beta).$$ Έτσι $(+2)-(+3)=(+2)+(-3)=-1$ και $(+4)-(-1)=(+4)+(+1)=+5$.
\end{kanonas}

Στους ακέραιους η αφαίρεση γίνεται πάντα: ο μειωτέος δεν χρειάζεται να είναι μεγαλύτερος από τον αφαιρετέο. Η διαφορά $\alpha-\beta$ δείχνει και τον δρόμο από το $\beta$ προς το $\alpha$ πάνω στον άξονα: από το $-3$ ως το $+5$ είναι $8$ βήματα δεξιά, και πράγματι $(+5)-(-3)=8$.

\tmimasmall{Απαλοιφή παρενθέσεων και αλγεβρικό άθροισμα}

Μια παράσταση με προσθέσεις και αφαιρέσεις γράφεται ως πρόσθεση ακεραίων: $12-15+8-20+3=12+(-15)+8+(-20)+3$. Τέτοια παράσταση ονομάζεται \textbf{αλγεβρικό άθροισμα} και μπορώ να τη λύσω με όποια σειρά θέλω: $(12+8+3)+(-15-20)=23-35=-12$. Για να σβήσω παρενθέσεις:
\begin{itemize}[leftmargin=1.2em,itemsep=0.3ex,topsep=2pt]
\item Όταν η παρένθεση έχει μπροστά της $+$ (ή δεν έχει πρόσημο), τη σβήνω μαζί με αυτό το $+$ και γράφω τους όρους της \textbf{με τα ίδια πρόσημα}: $(+5)+(-4)+(-2+6)=5-4-2+6$.
\item Όταν η παρένθεση έχει μπροστά της $-$, τη σβήνω μαζί με το $-$ και γράφω τους όρους της \textbf{με αντίθετα πρόσημα}: $-(+3)-(-4)-(-2+5)=-3+4+2-5$.
\end{itemize}

\begin{prosoxi}
\begin{itemize}[leftmargin=1.2em,itemsep=0.3ex,topsep=2pt]
\item Στους ακέραιους η πρόσθεση δεν σημαίνει πάντα αύξηση: $5+(-3)=2<5$. Και η αφαίρεση δεν σημαίνει πάντα μείωση: $3-(-2)=5>3$.
\item Στους ετερόσημους δεν προσθέτω τις απόλυτες τιμές: $(+9)+(-4)\neq13$, αλλά $9-4=5$.
\item Το πρόσημο του αθροίσματος ετεροσήμων είναι του μεγαλύτερου κατ' απόλυτη τιμή: $(-9)+4=-5$, όχι $+5$.
\item Η διαφορά δεν είναι αντιμεταθετική: $5-3=2$ αλλά $3-5=-2$.
\item Με το $-$ μπροστά από παρένθεση αλλάζω το πρόσημο \emph{όλων} των όρων της παρένθεσης, όχι μόνο του πρώτου: $10-(3-8+2)=10-3+8-2$.
\item Η διαφορά θερμοκρασιών $+5$ και $-3$ είναι $8$, όχι $2$. Αφαιρώ τον αριθμό μαζί με το πρόσημό του.
\end{itemize}
\end{prosoxi}

\tmima{Λυμένα παραδείγματα}

\paradeigma{Πρόσθεση στον άξονα}
\tixero Οι προσθέσεις $(-2)+5$, $3+(-7)$, $(-4)+(-3)$.\par
\tithelo Να τις κάνω στον άξονα.\par
\possindeo Ξεκινάω από τον πρώτο προσθετέο και κινούμαι: από το $-2$ πέντε δεξιά: $3$. Από το $3$ επτά αριστερά: $-4$. Από το $-4$ τρία αριστερά: $-7$.\par
\apantisi $3$, $-4$, $-7$.\par
\elegxos Με τους κανόνες: $5-2=3$, $-(7-3)=-4$, $-(4+3)=-7$.

\paradeigma{Με τους κανόνες}
\tixero Οι προσθέσεις $(-8)+(-5)$, $(+9)+(-4)$, $(-9)+(+4)$, $(-6)+(+6)$.\par
\tithelo Να υπολογίσω τα αθροίσματα.\par
\possindeo Ομόσημοι: $-(8+5)=-13$. Ετερόσημοι με μεγαλύτερη απόλυτη τιμή το $+9$: $+(9-4)=5$. Ετερόσημοι με μεγαλύτερη το $-9$: $-(9-4)=-5$. Αντίθετοι: $0$.\par
\apantisi $-13$, $5$, $-5$, $0$.\par
\elegxos Κάθε πρόσημο το βρίσκω από τον προσθετέο με τη μεγαλύτερη απόλυτη τιμή.

\paradeigma{Αφαιρέσεις}
\tixero Οι διαφορές $(+5)-(-3)$, $(+2)-(+3)$, $(-4)-(-6)$, $(-7)-(+5)$.\par
\tithelo Να τις υπολογίσω.\par
\possindeo Τις γράφω ως προσθέσεις του αντίθετου: $5+3$, $2+(-3)$, $(-4)+6$, $(-7)+(-5)$.\par
\apantisi $8$, $-1$, $2$, $-12$.\par
\elegxos Για την πρώτη: $(-3)+8=5$. Για την τελευταία: $(-12)+5=-7$.

\paradeigma{Το ταμείο της παρέας}
\tixero Το ταμείο έχει $+30$ €. Κινήσεις: $-12$, $+8$, $-35$, $+20$.\par
\tithelo Το υπόλοιπο.\par
\possindeo Θετικοί: $30+8+20=58$. Αρνητικοί: $-12-35=-47$. $58+(-47)=11$.\par
\apantisi $11$ €.\par
\elegxos Βήμα βήμα: $30,\ 18,\ 26,\ -9,\ 11$. Συμφωνεί.

\paradeigma{Παρενθέσεις με δύο τρόπους}
\tixero Η παράσταση $-(6-9+2)+(-5+1)-(4-10-3)$.\par
\tithelo Να την υπολογίσω.\par
\possindeo \textbf{1ος τρόπος:} κάνω πρώτα τις πράξεις μέσα στις παρενθέσεις: $6-9+2=-1$, $-5+1=-4$, $4-10-3=-9$. Άρα $-(-1)+(-4)-(-9)=1-4+9=6$. \textbf{2ος τρόπος:} σβήνω τις παρενθέσεις: $-6+9-2-5+1-4+10+3$. Θετικοί: $9+1+10+3=23$. Αρνητικοί: $-6-2-5-4=-17$. Άρα $23-17=6$.\par
\apantisi $6$.\par
\elegxos Και οι δύο τρόποι δίνουν το ίδιο αποτέλεσμα.

\paradeigma{Παραστάσεις με γράμματα}
\tixero Οι αριθμοί $\alpha=-2$, $\beta=+4$, $\gamma=-6$.\par
\tithelo Να υπολογίσω τις παραστάσεις $\alpha+\beta+\gamma$, $\beta-\alpha-\gamma$ και $\alpha-(\beta-\gamma)$.\par
\possindeo Βάζω τον αριθμό στη θέση του γράμματος, μέσα σε παρένθεση για να φαίνεται το πρόσημο. $\alpha+\beta+\gamma=(-2)+4+(-6)=-4$. $\beta-\alpha-\gamma=4-(-2)-(-6)=4+2+6=12$. Για την τρίτη, πρώτα $\beta-\gamma=4-(-6)=10$ και μετά $\alpha-10=(-2)-10=-12$.\par
\apantisi $-4$, $12$ και $-12$.\par
\elegxos Για την πρώτη: $(-2)+(-6)=-8$ και $-8+4=-4$.

\tmima{Ασκήσεις}

\zoni{Το κατάλαβα}{κανόνες, ιδιότητες, παρενθέσεις}

\ex Στον άξονα βρες τα αθροίσματα: α) $(-3)+5$ \quad β) $2+(-6)$ \quad γ) $(-5)+(-2)$ \quad δ) $7+(-7)$ \quad ε) $0+(-4)$.

\ex Υπολόγισε: α) $(-8)+(-5)$ \quad β) $(+9)+(-4)$ \quad γ) $(-9)+(+4)$ \quad δ) $(+7)+(+5)$ \quad ε) $(-12)+0$ \quad στ) $(-15)+(+15)$ \quad ζ) $(+20)+(-35)$ \quad η) $(-17)+(+6)$.

\ex Υπολόγισε τις διαφορές: α) $(+3)-(+1)$ \quad β) $(+10)-(+20)$ \quad γ) $(+15)-(-5)$ \quad δ) $(-6)-(+2)$ \quad ε) $(-11)-(-6)$ \quad στ) $(+8)-(-8)$ \quad ζ) $(-14)-(-14)$ \quad η) $(-25)-(+8)$ \quad θ) $(-9)-(-17)$.

\ex Συμπλήρωσε τον αριθμό που λείπει ώστε το άθροισμα να είναι $0$: α) $(-7)+\kouti$ \quad β) $(+13)+\kouti$ \quad γ) $\kouti+(-4)$ \quad δ) $(-21)+\kouti$.

\ex Για κάθε ζεύγος $(\alpha,\beta)$ βρες τα $\alpha+\beta$, $\alpha-\beta$ και $\beta-\alpha$: α) $(+6,\,+2)$ \quad β) $(-12,\,-7)$ \quad γ) $(+9,\,-14)$ \quad δ) $(-2,\,+5)$. Τι παρατηρείς για τους αριθμούς $\alpha-\beta$ και $\beta-\alpha$;

\ex Υπολόγισε: α) $(-7)+12+(-5)+9+(-3)$ \quad β) $(-15)+4+(-25)+11$ \quad γ) $10+(-3)+(-4)+(-5)$.

\ex Σβήσε τις παρενθέσεις και υπολόγισε: α) $5+(-4)+(-2+6)$ \quad β) $-(+3)-(-4)-(-2+5)$ \quad γ) $(14-9+1)-(13-9)+3$ \quad δ) $12-(5-8+3)-(-4+9)$.

\ex Σωστό ή λάθος; α) Αν το άθροισμα δύο ακεραίων είναι θετικό, τότε και οι δύο είναι θετικοί. \quad β) Αν το άθροισμα δύο ακεραίων είναι αρνητικό, τότε και οι δύο είναι αρνητικοί. \quad γ) Η πρόσθεση στους ακέραιους σημαίνει πάντα αύξηση. \quad δ) Η αφαίρεση στους ακέραιους σημαίνει πάντα μείωση. \quad ε) Αν η διαφορά δύο ακεραίων είναι $0$, οι αριθμοί είναι ίσοι. \quad στ) $-(-1)=1$. \quad ζ) Το άθροισμα δύο αντίθετων αριθμών είναι $0$. \quad η) $5-7=7-5$.

\zoni{Το χρησιμοποιώ}{θερμόμετρα, ανελκυστήρες, ταμεία, παιχνίδια}

\ex \textbf{Η θερμοκρασία.} Το πρωί στο Μέτσοβο είναι $-6$°, το μεσημέρι ανεβαίνει $9$ βαθμούς και το βράδυ πέφτει $5$ βαθμούς. Πόση είναι η θερμοκρασία το μεσημέρι και το βράδυ;

\ex \textbf{Ο ανελκυστήρας.} Ο Ηλίας είναι στον όροφο $-2$, ανεβαίνει $5$ ορόφους και μετά κατεβαίνει $4$. Σε ποιον όροφο φτάνει;

\ex \textbf{Το παιχνίδι.} Οι πόντοι της Πολυξένης στους τέσσερις γύρους είναι $+8$, $-5$, $-7$, $+12$. Πόσους έχει συνολικά;

\ex \textbf{Ο λογαριασμός.} Η Ρήνη έχει υπόλοιπο $-45$ € στον λογαριασμό, καταθέτει $80$ € και μετά αναλαμβάνει $20$ €. Πόσα έχει;

\ex \textbf{Θερμική διαφορά.} Στο Μέτσοβο η μέγιστη θερμοκρασία ήταν $+4$° και η ελάχιστη $-6$°. Στη Φλώρινα η μέγιστη ήταν $-1$° και η ελάχιστη $-9$°. Σε ποια πόλη ήταν μεγαλύτερη η διαφορά μέγιστης και ελάχιστης;

\ex \textbf{Από την περιγραφή στην πράξη.} Γράψε μια παράσταση για καθεμιά από τις καταστάσεις και υπολόγισέ την. α) Ένα σημείο ξεκινά από το $-4$, κινείται $6$ μονάδες δεξιά και μετά $3$ αριστερά. β) Ο Βασίλης βρίσκεται στα $15$ μ. υψόμετρο, κατεβαίνει $25$ μ. και μετά ανεβαίνει $6$ μ. γ) Ο λογαριασμός της Φανής έχει $120$ €, γίνεται κατάθεση $40$ € και ανάληψη $90$ €.

\zoni{Το σκέφτομαι}{γιατί ισχύουν οι κανόνες}

\ex Πότε το άθροισμα δύο ακεραίων είναι $0$;

\ex Αν οι $\alpha,\,\beta$ είναι ετερόσημοι, ποιο πρόσημο έχει το άθροισμά τους;

\ex Είναι πάντα το άθροισμα μεγαλύτερο από κάθε προσθετέο; Δώσε αντιπαράδειγμα.

\ex Ο Στέλιος γράφει $5-(-3)=2$, γιατί «πέντε μείον τρία κάνει δύο». Ποιο είναι το λάθος του και ποιο το σωστό αποτέλεσμα;

\zoni{Για το σπίτι}{ο προϋπολογισμός σου}

\ex Κράτησε για μια εβδομάδα τα έσοδα και τα έξοδα του χαρτζιλικιού σου ως ακέραιους και υπολόγισε το τελικό υπόλοιπο.

\ex Φτιάξε ένα παιχνίδι με κάρτες $+$ και $-$ και παίξε με έναν φίλο τρεις γύρους. Ποιος έχει το μεγαλύτερο άθροισμα;
